(b) It is clear from the definition (4) that
\[
p_n(x/y)_{x_1=t,\,y_1=-t}=p_n(x/y)_{x_1=y_1=0}.
\]
Since $\omega_y$ and the substitutions $x_1=t$, $y_1=-t$ and $x_1=y_1=0$ are homomorphisms (into suitable rings), it follows that $g(x,y)$ satisfies (6) whenever $g\in\operatorname{im}(\omega_y)$.

Conversely, suppose that $g(x,y)\in\Lambda(x)\otimes\Lambda(y)$ and that $g(x,y)$ satisfies (6). Thus by hypothesis
\[
g((t,x),(-t,y))=g(x,y), \tag{25}
\]
where $(t,x)=(t,x_1,x_2,\ldots)$ and similarly for $(t,y)$. Define
\[
\begin{aligned}
u_i&=\frac{1}{2}\Big(p_i(x)+(-1)^i p_i(y)\Big)\\
v_i&=\frac{1}{2}\Big(p_i(x)+(-1)^{i-1}p_i(y)\Big).
\end{aligned}
\]
It is clear that $g(x,y)$ can be written uniquely as a polynomial $P(u_1,u_2,\ldots;v_1,v_2,\ldots)$ in the $u_i$'s and $v_i$'s, and we want to show that no $u_i$ appears. Now
\[
g((t,x),(-t,y))=P(u_1+t,u_2+t^2,\ldots;v_1,v_2,\ldots). \tag{26}
\]
By considering the coefficient of $t$ in (26) and using (25), we see that $u_1$ does not appear in $P$. Now by considering the coefficient of $t^2$ is follows that $u_2$ does not appear in $P$, etc. This proof is due to John Stembridge (unpublished).
